% !TEX program = xelatex
% 数学物理方法1 · 第五章 · 选题 8（课件 p49--p50 例 1）卷积反解 · 解答卡
\documentclass[10pt]{ctexart}
%=====================================================================
% theme-violet.tex —— 讲义精讲 · 紫花主题
% 用法：\documentclass[11pt]{ctexart}
%       \input{assets/theme-violet}
%       \begin{document} ... \end{document}
% 只用 ctexart 自带能力 + xcolor/enumitem/fancyhdr/pifont/amsmath，
% 不依赖 tikz/tcolorbox，保证 MiKTeX 可直接编译。
%=====================================================================

\usepackage[a4paper,left=2.0cm,right=2.0cm,top=2.1cm,bottom=2.3cm,%
            headsep=12pt,footskip=22pt]{geometry}
\usepackage{amsmath,amssymb}
\usepackage{needspace}
\usepackage{xcolor}
\usepackage{array}
\usepackage{enumitem}
\usepackage{pifont}
\usepackage{fancyhdr}
\usepackage[hidelinks]{hyperref}

%---------------------------------------------------------------- 配色
\definecolor{violet}{RGB}{107,78,143}      % 主色：紫
\definecolor{violetdeep}{RGB}{63,44,92}    % 深紫：标题/强调
\definecolor{petal}{RGB}{183,159,208}      % 花瓣：装饰线、花饰
\definecolor{bandviolet}{RGB}{247,243,251} % 浅紫底：档案/提示框
\definecolor{bandcream}{RGB}{252,249,242}  % 米白底：易错/注意框
\definecolor{rulepurple}{RGB}{217,207,230} % 分隔线
\definecolor{inkgray}{RGB}{92,96,110}      % 次要文字

%---------------------------------------------------------------- 字体
% 有 Noto CJK 用 Noto，没有就退回 ctex 默认（Windows 下为宋体/黑体）
\IfFontExistsTF{Noto Serif CJK SC}{\setCJKmainfont{Noto Serif CJK SC}}{}
\IfFontExistsTF{Noto Sans CJK SC}{\setCJKsansfont{Noto Sans CJK SC}}{}

%---------------------------------------------------------------- 版式
\setlist{nosep,leftmargin=1.9em}
\linespread{1.28}
\setlength{\parskip}{3.5pt}
\setlength{\parindent}{0pt}
\setlength{\tabcolsep}{5pt}
\renewcommand{\arraystretch}{1.28}

%---------------------------------------------------------------- 花饰
% \ding{95} = 五瓣小花，\ding{96} = 星形花；都来自 pifont，必定存在
\newcommand{\orn}{\textcolor{petal}{\ding{95}}}          % 行内小标
\newcommand{\ornstar}{\textcolor{petal}{\ding{96}}}      % 星形花
\newcommand{\flowerrule}{%                                % 居中花饰分隔线
  \par\vspace{7pt}%
  \noindent\makebox[\linewidth]{\color{petal}%
    \rule{0.20\linewidth}{0.6pt}\hspace{8pt}{\small\ding{95}}\hspace{4pt}%
    \raisebox{1.1pt}{\tiny\ding{95}}\hspace{4pt}{\small\ding{95}}\hspace{8pt}%
    \rule{0.20\linewidth}{0.6pt}}%
  \par\vspace{7pt}}

%---------------------------------------------------------------- 标题与节头
\newcommand{\doctitle}[3]{%   #1 主标题  #2 副标题  #3 元信息（可用 \\ 分行）
  \begin{center}
    {\color{petal}\small \orn\enspace #2\enspace \ornstar}\\[7pt]
    {\LARGE\bfseries\color{violetdeep} #1}\\[10pt]
    {\color{petal}\rule{0.55\linewidth}{1pt}}\\[4pt]
    {\small\color{inkgray} #3}
  \end{center}\vspace{2pt}}

\newcommand{\hdr}[2]{%   #1 章节号（如 §0）  #2 标题
  \par\vspace{12pt}%
  \noindent\colorbox{violetdeep}{\color{white}\bfseries\small\;#1\;}\hspace{7pt}%
  {\color{violetdeep}\large\bfseries #2}\par\nobreak
  \vspace{3pt}{\color{petal}\hrule height 1.2pt}\vspace{6pt}}

\newcommand{\subhdr}[1]{%
  \par\vspace{6pt}%
  \noindent{\color{violet}\bfseries #1}\par\vspace{1pt}}

%---------------------------------------------------------------- 提示框
\newcommand{\callout}[1]{% 浅紫底：讲义档案、读法说明
  \par\medskip\noindent
  \fcolorbox{rulepurple}{bandviolet}{\parbox{\dimexpr\linewidth-2\fboxsep-2\fboxrule}{\vspace{2pt}#1\vspace{2pt}}}%
  \par\medskip}

\newcommand{\warnbox}[1]{% 米白底：易错、边界条件
  \par\medskip\noindent
  \fcolorbox{rulepurple}{bandcream}{\parbox{\dimexpr\linewidth-2\fboxsep-2\fboxrule}{\vspace{2pt}#1\vspace{2pt}}}%
  \par\medskip}

\newcommand{\gapitem}[2]{% #1 缺口类型（紫） #2 内容
  \par\vspace{2pt}\noindent
  \hangindent=1.4em\hangafter=1
  {\color{petal}\small\ding{96}}\enspace{\bfseries\color{violetdeep}#1}\enspace #2\par}

% 题干块（\qstem）：每道例题／作业的「详解」开头必须先复述题干与出处。
% 起因：2026-10-25 第五章讲义把七道作业的题面集中放在一页、详解另起一页，
% 读者翻到「第 4 题」时看不到条件，等于一道没有题干的题目。
% 规则：详解块内自带题干，读者不翻页就能读懂题在问什么。
\newcommand{\qstem}[1]{%
  \par\medskip\noindent
  \fcolorbox{violet}{bandviolet}{%
    \parbox{\dimexpr\linewidth-2\fboxsep-2\fboxrule}{%
      \vspace{2pt}%
      {\small\bfseries\color{violet}题干}\hspace{0.8em}#1%
      \vspace{2pt}}}%
  \par\medskip}

% v2: 保护题干与开头，长题可跨页，续页补上下文。
\newcommand{\qbreak}{\needspace{6\baselineskip}}
% 独立练习：ID、来源、题干、考点、答案标签（通常 ans:ID）。
\newcommand{\practice}[5]{%
  \qbreak\subhdr{练习 #1}\qstem{（#2）#3}%
  \textbf{考点}：#4\par
  详解见附录第~\pageref{#5}~页。\par\medskip}
% 附录开头：ID、来源、条件摘要；自动产生 ans:ID 标签。
\newcommand{\answerstart}[3]{%
  \qbreak\subhdr{#1\quad 详解}\label{ans:#1}\qstem{（#2）#3}}
\newcommand{\answercontinue}[2]{%
  \needspace{5\baselineskip}\subhdr{#1\quad 续解}%
  {\small\color{inkgray}本页接续：#2}\par\medskip}
% 延迟提供默认值，让旧讲义可在导言区保留自己的 \newcommand{\step}。
\AtBeginDocument{\providecommand{\step}[1]{\needspace{3\baselineskip}\par\medskip\textbf{#1}\par}}

% 前提旁注：只放在「需要理解的那一步」后面，给陈述 + 本题代入，不重证。
\newcommand{\aside}[1]{%
  \par\vspace{1pt}\noindent
  \hangindent=1.6em\hangafter=1
  {\small\color{inkgray}%
    \hspace*{0.6em}{\color{petal}$\triangleright$}\hspace{2pt}#1}%
  \par\vspace{2pt}}

%---------------------------------------------------------------- 表格
% 本机 MiKTeX 的 colortbl 与 array 版本不匹配（\rowcolor 会报错），
% 故表格统一用「上粗紫线 + 表头 + 中细线」的三线式，不依赖 colortbl。
\newcommand{\tabtop}{\par\vspace{3pt}{\color{violet}\hrule height 1.2pt}\vspace{1pt}}
\newcommand{\tabmid}{\par\vspace{1pt}{\color{petal}\hrule height 0.5pt}\vspace{1pt}}
\newcommand{\tabbot}{\par\vspace{3pt}{\color{rulepurple}\hrule height 0.5pt}\vspace{8pt}}
\newcommand{\thead}[1]{{\bfseries\color{violetdeep}#1}}

%---------------------------------------------------------------- 页眉页脚
\pagestyle{fancy}
\fancyhf{}
\renewcommand{\headrulewidth}{0.4pt}
\renewcommand{\headrule}{\hbox to\headwidth{\color{rulepurple}\leaders\hrule height \headrulewidth\hfill}}
\fancyhead[L]{\small\color{inkgray}\songti \doctag}
\fancyhead[R]{\small\color{inkgray}\songti \coursetag}
\fancyfoot[C]{\small\color{inkgray}\orn\enspace{\color{violet}\thepage}\enspace\orn}

% 每篇讲义在导言区设这两个值
\newcommand{\doctag}{}
\newcommand{\coursetag}{}
\newcommand{\settags}[2]{\renewcommand{\doctag}{#1}\renewcommand{\coursetag}{#2}}

%---------------------------------------------------------------- 数学记号
\newcommand{\vect}[1]{\boldsymbol{#1}}
\newcommand{\dd}{\,\mathrm{d}}
\newcommand{\ddx}{\frac{\mathrm{d}}{\mathrm{d}x}}
\addtolength{\jot}{2pt}



\settags{第五章 · 选题 8}{数学物理方法1}

\geometry{left=1.8cm,right=1.8cm,top=1.9cm,bottom=2.0cm,headsep=10pt,footskip=18pt}
\linespread{1.22}
\setlength{\headheight}{22pt}
\setlength{\parskip}{0pt}
\setlength{\abovedisplayskip}{3pt plus 1pt minus 1pt}
\setlength{\belowdisplayskip}{3pt plus 1pt minus 1pt}
\setlength{\abovedisplayshortskip}{2pt}
\setlength{\belowdisplayshortskip}{2pt}
\newcommand{\step}[1]{\par\smallskip\noindent{\color{violetdeep}\bfseries #1}\hspace{0.35em}}
\newcommand{\ok}{\textcolor{violet}{\ensuremath{\checkmark}}}

% 紧凑节头：留白减半，并防止节头孤零零留在页底
\newcommand{\chdr}[2]{%
  \needspace{4\baselineskip}%
  \par\vspace{3pt}%
  \noindent\colorbox{violetdeep}{\color{white}\bfseries\small\;#1\;}\hspace{7pt}%
  {\color{violetdeep}\bfseries #2}\par\nobreak
  \vspace{1pt}{\color{petal}\hrule height 1pt}\vspace{3pt}}
\let\hdr\chdr

\begin{document}

\doctitle{卷积定理反解：把 $|k|$ 请回来}
{未知函数被卷积“藏”起来，除一下，再逆变换}
{依据课件 p49--p50 例 1（第五章习题册第 8 题）\quad 2026-11-06\quad 数学物理方法\,1}

\qstem{已知 $\displaystyle\int_{-\infty}^{\infty}\frac{f(\eta)}{(x-\eta)^{2}+a^{2}}\dd\eta=\frac{1}{x^{2}+b^{2}}$，且 $0<a<b$，求 $f(x)$。}

\hdr{§1}{认题：左端就是一个卷积}

$f$ 的自变量是 $\eta$，核的自变量是 $x-\eta$，只差一个平移——这正是卷积的定义形状。记 $f_{2}(x)=1/(x^{2}+a^{2})$，原式即 $f*f_{2}=1/(x^{2}+b^{2})$：卷积定理把“求函数”换成了“求像函数”。

\hdr{§2}{两件工具：本章约定 A 与泊松核}

本章 $F$ 的定义式带 $\dfrac{1}{2\pi}$，卷积定理因此还回来一个 $2\pi$（课件 p46，\textbf{两条不对称，必须配套用}）：
\[
F(k)=\frac{1}{2\pi}\int_{-\infty}^{\infty}f(x)e^{-ikx}\dd x,\qquad f(x)=\int_{-\infty}^{\infty}F(k)e^{ikx}\dd k,
\qquad f_{1}*f_{2}\rightleftharpoons 2\pi F_{1}F_{2}.
\]
核的变换用留数算。因 $\left|e^{-ikz}\right|=e^{ky}$（$z=x+iy$），闭合方向被 $k$ 的符号定死：
\[
\mathcal F\!\left[\frac{1}{x^{2}+c^{2}}\right]
=\frac{1}{2\pi}\cdot
\begin{cases}
+2\pi i\operatorname{Res}_{z=ci} & k<0\ \text{（上半平面，逆时针）}\\[2pt]
-2\pi i\operatorname{Res}_{z=-ci} & k>0\ \text{（下半平面，顺时针）}
\end{cases}
=\frac{e^{-|k|c}}{2c}.
\]
两条分支给出同一条结果——\textbf{绝对值不是“补”上去的，是分支算完的结果}。

\hdr{§3}{代入卷积定理，把像函数解出来}

对 $f*f_{2}=\dfrac{1}{x^{2}+b^{2}}$ 两边取变换，左端用卷积定理：
\[
2\pi\,\mathcal F[f]\cdot\frac{e^{-|k|a}}{2a}=\frac{e^{-|k|b}}{2b}
\quad\Longrightarrow\quad
\mathcal F[f]=\frac{1}{2\pi}\cdot\frac{e^{-|k|b}}{2b}\cdot\frac{2a}{e^{-|k|a}}
=\frac{a}{2\pi b}\,e^{-|k|(b-a)} .
\]
分母对每个 $k$ 都非零，所以能“除”过去；而 $b>a$ 的作用，正是让 $e^{-|k|(b-a)}$ 在 $|k|\to\infty$ 时衰减。

\hdr{§4}{反演：认出标准形，不必硬算积分}

令 $c=b-a>0$，把 $\dfrac{a}{2\pi b}e^{-|k|c}$ 凑成 §2 的标准形 $\dfrac{e^{-|k|c}}{2c}$：
\[
\frac{a}{2\pi b}\,e^{-|k|c}=\frac{a(b-a)}{\pi b}\cdot\frac{e^{-|k|c}}{2(b-a)}
\quad\Longrightarrow\quad
\boxed{\;f(x)=\frac{a(b-a)}{\pi b}\cdot\frac{1}{x^{2}+(b-a)^{2}}\;}
\]
\textbf{自检：}$\int f\dd x=a/b<1$ \ok{}；$b\to a^{+}$ 时 $f\to\delta(x)$、原式退化为恒等式 \ok{}；$a=1,b=2,x=0$ 时左端 $=1/4=1/(0+2^{2})$ \ok{}（$x=0.5,1,2.5$ 亦已数值验过）。

\hdr{§5}{对照你草稿上的三处卡点}

\gapitem{卡点一 · 定义式的 $\frac{1}{2\pi}$}{草稿把 $\mathcal F\!\left[\frac{1}{x^{2}+a^{2}}\right]$ 写成 $\frac{\pi}{a}e^{ka}$，正确值是 $\frac{e^{-ka}}{2a}$——大 $2\pi$ 倍。算完 $2\pi i\cdot\operatorname{Res}$ 后，漏掉了定义式前面那个 $\frac{1}{2\pi}$。}
\gapitem{卡点二 · 丢掉的 $|k|$}{写成 $e^{ka}$、$e^{kb}$ 而不带绝对值，等于默认 $k$ 只有一种符号。逆变换要让 $k$ 跑遍 $(-\infty,\infty)$，反向的一半 $e^{k(b-a)}$ 指数发散——你草稿卡在 $\int e^{ikx+k(b-a)}\dd k$ 算不下去，就是这个原因。}
\gapitem{卡点三 · 硬算逆变换}{不必对 $e^{-|k|(b-a)+ikx}$ 直接积分：标准形 $\frac{e^{-|k|c}}{2c}$ 就是 $\frac{1}{x^{2}+c^{2}}$ 的像，认出即报。}

\callout{\textbf{一句话带走}：卷积定理只做两件事——卷积变乘积、求函数变求像函数；工作量全在\emph{算好那一个核的变换}，而它一定带 $|k|$。}

\end{document}
